\chapter{I sensori di presenza: principi e stato dell'arte}
\label{cap:sensori}

\begin{quote}\small
Questo capitolo risponde all'obiettivo O1: che cosa sono i sensori di presenza,
come funzionano, chi li usa e quali dati producono. La trattazione è
volutamente asimmetrica --- più spazio a PIR e mmWave, che sono i sensori
effettivamente misurati nel \cref{cap:confronto} --- e si chiude con il
posizionamento comparativo che giustifica le scelte del resto della tesi.
\end{quote}

\section{Il sensore PIR}
\label{sec:pir}

\subsection{Principio fisico: un rilevatore di variazione, non di presenza}
\label{sec:pir-principio}

PIR è l'acronimo di \emph{Passive InfraRed}: il sensore non emette nulla, si limita
a osservare la radiazione infrarossa che lo raggiunge. Il cuore del componente è un
cristallo \textbf{piroelettrico}, che genera una piccola carica elettrica quando la
radiazione incidente \emph{cambia}. Il corpo umano, con una temperatura superficiale
di circa 34\,\celsius, emette con picco intorno a 9,4~\textmu m; il sensore è
filtrato con una finestra ottica di 5--14~\textmu m proprio per privilegiare quella
banda~\cite{adafruitPIR}.

Il punto essenziale per questa tesi è che il cristallo risponde alla
\textbf{derivata} del flusso infrarosso, non al suo valore assoluto.
Costruttivamente il sensore contiene \textbf{due elementi piroelettrici in
configurazione differenziale}, collegati in serie con polarità opposte, e da questo
seguono tre conseguenze:

\begin{itemize}[leftmargin=1.4em,itemsep=2pt]
  \item una persona che si muove investe i due elementi in \emph{tempi diversi}: il
        segnale differenziale è non nullo e si ha rilevamento;
  \item una persona ferma illumina entrambi gli elementi con lo stesso flusso
        costante: la differenza è nulla e il sensore è \textbf{fisicamente cieco};
  \item una variazione globale della scena (il sole che scalda la stanza) colpisce
        entrambi gli elementi allo stesso modo e si cancella: da qui la buona
        immunità del PIR ai falsi positivi ambientali lenti.
\end{itemize}

È importante inquadrare correttamente questo comportamento. Il risultato
sperimentale del \cref{sec:persona-ferma} --- falsi negativi prossimi al
100\,\% su persona ferma --- \textbf{non è il fallimento di un componente
economico}: è il funzionamento nominale del principio. Nessun PIR, a nessun prezzo,
può rilevare un corpo immobile. Questa è la frase su cui poggia l'intero confronto
con il radar.

\subsection{La lente di Fresnel: da due elementi a decine di zone}
\label{sec:fresnel}

Due soli elementi coprirebbero un campo visivo strettissimo. La cupola bianca che
si vede su tutti i moduli PIR è una \textbf{lente di Fresnel multi-segmento}: ogni
segmento focalizza sul cristallo una porzione differente della scena, creando
decine di \emph{zone di rilevamento} a ventaglio, separate da zone cieche.

Questo dettaglio costruttivo ha due conseguenze misurabili, entrambe verificate
sperimentalmente nel \cref{cap:confronto}:

\begin{enumerate}[leftmargin=1.8em,itemsep=2pt]
  \item il movimento \textbf{trasversale}, che attraversa le zone, genera molte
        transizioni di flusso ed è rilevato bene; il movimento \textbf{radiale},
        diretto verso il sensore, resta più a lungo nella stessa zona ed è rilevato
        male. Il radar FMCW ha la sensibilità esattamente opposta, perché misura
        variazioni di distanza;
  \item la portata dichiarata (3--7~m) vale per il movimento trasversale di tutto
        il corpo. Un movimento \emph{sul posto} non attraversa le zone e cade nel
        caso sfavorevole: è il meccanismo che spiega il risultato, a prima vista
        sorprendente, del \cref{sec:pir-movimento}.
\end{enumerate}

\subsection{I dati prodotti: un bit, e come va interpretato}
\label{sec:pir-dati}

L'uscita di un modulo PIR è \textbf{un solo bit digitale}. Tutta l'elaborazione è
nell'elettronica di condizionamento, tipicamente il circuito integrato
BISS0001~\cite{biss0001}, che introduce i meccanismi riassunti nella
\cref{tab:pir-meccanismi}.

\begin{table*}[htbp]
\centering
\caption{Meccanismi di condizionamento del segnale PIR e loro effetto sul dato
prodotto.}
\label{tab:pir-meccanismi}
\small
\begin{tabularx}{\linewidth}{@{}lX@{}}
\toprule
\textbf{Meccanismo} & \textbf{Effetto sul dato} \\
\midrule
Soglia sul segnale amplificato & sensibilità/portata (trimmer) \\
\textbf{Tempo di ritenuta} (hold) & l'uscita resta alta $N$ secondi dopo l'ultimo
  trigger (trimmer) \\
Modalità trigger L (single) & l'uscita torna bassa alla fine della ritenuta anche
  se il movimento continua \\
Modalità trigger H (repeat) & ogni nuovo movimento riarma la ritenuta \\
Tempo di stabilizzazione & 30--60~s all'accensione, uscita non attendibile \\
Block time & 2,5~s di cecità strutturale dopo ogni rilascio \\
\bottomrule
\end{tabularx}
\end{table*}

Da questa tabella derivano direttamente tre scelte del protocollo sperimentale
(\cref{sec:protocollo}): il jumper in modalità \textbf{H}, il trimmer di ritenuta
al \textbf{minimo} e lo scarto del primo minuto di ogni acquisizione. Se la
ritenuta non fosse al minimo, la ``latenza di rilascio'' misurata sarebbe una
proprietà del potenziometro, non del sensore.

Il confronto informativo con il mmWave si riassume in una riga: il PIR produce
\textbf{1 bit filtrato da isteresi}, il LD2410B produce circa \textbf{venti valori
numerici} (stato, due distanze, due energie di target, diciotto energie per-gate) a
ogni frame. La tesi quantifica quanto questo divario informativo si traduca in
capacità di rilevamento.

\subsection{Sensibilità alla temperatura ambiente}
\label{sec:pir-temperatura}

Il segnale utile del PIR è il \emph{contrasto termico} fra corpo e sfondo. A
20\,\celsius\ di ambiente il contrasto è di circa 14\,\celsius; a
30--32\,\celsius\ --- un'aula assolata d'estate --- scende a 2--4\,\celsius, con
segnale più debole, portata ridotta e più falsi negativi. Sopra i 35\,\celsius\ il
contrasto può addirittura invertirsi.

Per questo motivo ogni sessione sperimentale registra la temperatura della stanza
(\cref{sec:registro}). Le acquisizioni di questa tesi sono state svolte fra 25 e
29\,\celsius: una condizione \emph{sfavorevole} al PIR, che va dichiarata
esplicitamente nella discussione dei risultati (\cref{sec:limiti}).

Altre sorgenti note di falsi positivi PIR, tenute fuori dal setup: correnti d'aria
calda o fredda, animali, sole diretto attraverso i vetri (il vetro blocca l'IR del
corpo, ma non il riscaldamento solare degli oggetti interni).

\subsection{Chi lo usa e perché domina}
\label{sec:pir-usi}

Il PIR è il sensore di presenza più diffuso al mondo: luci automatiche, antifurti,
videocitofoni, termostati. Le ragioni sono tre numeri difficili da battere:
circa \textbf{50~\textmu A} di corrente di riposo, \textbf{1--3~\euro} di costo e
\textbf{zero elaborazione} richiesta al microcontrollore, cui si aggiunge l'assenza
totale di emissioni --- essendo un sensore passivo --- che semplifica
certificazioni e conformità in materia di privacy.

Nel progetto UPRISE il DIPME-DEVICE monta un PIR esattamente per questi motivi: un
nodo deve poter vegliare per anni a batteria. Il limite emerge solo nello scenario
di emergenza, ed è il punto di partenza di questa tesi.

\section{Il radar mmWave a 24~GHz}
\label{sec:mmwave}

\subsection{Principio: FMCW e gate di distanza}

I moduli in esame sono radar \emph{FMCW} (\emph{Frequency Modulated Continuous
Wave}) operanti nella banda ISM 24--24,25~GHz, come dichiarato dai rispettivi
manuali ufficiali: il \S7 del manuale del LD2410B riporta
\emph{operating frequency 24GHz\,\textasciitilde\,24.25GHz}~\cite{hilinkManualV104}
e la Tabella~2-1 del manuale del LD2420 indica la stessa banda con una banda di
sweep massima di 0,25~GHz~\cite{hilinkDriveLD2420}. Si tratta di
un'allocazione SRD/ISM: un modulo che dichiari conformità FCC e CE deve
necessariamente restare in quei 250~MHz, il che rende il dato verificabile
indipendentemente dalla scheda tecnica del produttore.\footnote{Una precisazione
terminologica: l'espressione \emph{mmWave}, di uso corrente in letteratura e
nella documentazione commerciale, è in senso stretto impropria a questa
frequenza. La banda delle onde millimetriche (EHF) va da 30 a 300~GHz, cioè
lunghezze d'onda fra 10 e 1~mm; a 24,125~GHz la lunghezza d'onda è
$\lambda = c/f \approx 12{,}4$~mm, dunque ancora in banda SHF. In questo lavoro
si mantiene la dizione \emph{mmWave} perché è quella adottata dai produttori e
dalla letteratura citata.}

Il principio è il seguente: il
trasmettitore emette un segnale continuo la cui frequenza cresce linearmente nel
tempo (\emph{chirp}); l'eco riflessa dal bersaglio torna con un ritardo
proporzionale alla distanza, e poiché la frequenza istantanea è cambiata nel
frattempo, il battimento fra segnale trasmesso e ricevuto ha una frequenza
proporzionale alla distanza del bersaglio. Un'analisi in frequenza del segnale di
battimento produce quindi direttamente un profilo distanza--intensità.

I moduli consumer non espongono il segnale grezzo, ma una sua versione elaborata e
discretizzata in \textbf{gate di distanza}: intervalli contigui di profondità, nove
per il LD2410B (0,75~m ciascuno, per un massimo di 6,75~m) e quindici per il
LD2420 (0,70~m ciascuno). Per ogni gate il modulo riporta un valore di
\textbf{energia riflessa} normalizzato nell'intervallo 0--100, separatamente per la
componente ``in movimento'' e per quella ``stazionaria''. Questo insieme di valori
--- accessibile solo attivando l'\emph{engineering mode}, \cref{sec:engineering} ---
è la materia prima di tutta l'analisi quantitativa di questa tesi.

\subsection{Come si rileva un corpo immobile}

Un radar FMCW puro misura distanza, non movimento: un corpo fermo produce comunque
un'eco. In pratica i moduli commerciali separano due contributi:

\begin{itemize}[leftmargin=1.4em,itemsep=2pt]
  \item la componente \textbf{moving}, associata a variazioni rapide della fase
        dell'eco (spostamento del bersaglio);
  \item la componente \textbf{stationary} (o \emph{motionless}), associata a
        un'eco stabile in distanza ma con fluttuazioni residue.
\end{itemize}

Una persona ferma non è mai realmente ferma: il torace si muove di alcuni
millimetri a ogni atto respiratorio, e a 24~GHz (lunghezza d'onda di circa
12,5~mm) uno spostamento millimetrico produce una variazione di fase perfettamente
misurabile. È questa la ragione fisica per cui il radar vede ciò che il PIR non può
vedere, ed è anche la base del rilevamento del respiro
(\cref{sec:respiro-teoria}).

\subsection{Il respiro come segnale}
\label{sec:respiro-teoria}

Il respiro di un adulto a riposo ha frequenza compresa fra circa 0,1 e 0,5~Hz
(6--30 atti al minuto). L'oscillazione toracica modula l'energia riflessa dal gate
in cui si trova la persona: campionando quel valore nel tempo e applicando una
trasformata di Fourier, il picco spettrale in banda 0,1--0,5~Hz corrisponde al
ritmo respiratorio.

Questa strada è ampiamente documentata in letteratura, con hardware sia
professionale sia consumer. Sono stati esaminati in particolare i lavori
sul monitoraggio non a contatto dei parametri vitali con radar FMCW, che
dimostrano il rilevamento di eventi apnoici e confrontano architetture a onda
continua e FMCW~\cite{upm2024vitals}, l'estrazione congiunta di frequenza
respiratoria e cardiaca con reti neurali~\cite{uts2024vitals}, il confronto fra
sensori a 24 e a 60~GHz (dove il 24~GHz privilegia la portata e il 60~GHz la
risoluzione)~\cite{theseus2024}, il tracciamento multi-bersaglio~\cite{uts2023mot}
e le tecniche specifiche di monitoraggio respiratorio con
mmWave~\cite{mdpi2024mmwaverm}.

L'aspettativa realistica con hardware consumer, che i risultati del
\cref{sec:respiro} confermano, è la seguente: rilevare la \emph{presenza} di
respiro --- cioè distinguere una persona ferma e viva da una stanza vuota --- è
affidabile; misurare la frequenza esatta richiede post-elaborazione e condizioni
controllate.

\subsection{Chi li usa}

I radar mmWave di fascia consumer si sono diffusi rapidamente nella domotica
(interruttori di presenza, illuminazione che non si spegne mentre si legge fermi
sul divano), nell'automazione degli edifici per il conteggio di occupazione, nel
monitoraggio del sonno e nei sistemi di rilevamento cadute per l'assistenza agli
anziani. Il supporto nelle piattaforme aperte è ormai maturo: le integrazioni
ESPHome per LD2410 e LD2420 sono documentate e ampiamente
usate~\cite{esphomeLD2410,esphomeLD2420,espboards}, e la letteratura tecnica sulle
applicazioni di conteggio dell'occupazione negli edifici
intelligenti~\cite{mmwaveOccupancy} documenta gli ordini di grandezza attesi. Sono
fonti utili per il
comportamento pratico dei moduli oltre a quanto scritto nei manuali.

\section{Altre tecnologie di rilevamento di presenza}
\label{sec:altre-tecnologie}

\subsection{UWB (Ultra-WideBand)}
\label{sec:uwb}

L'UWB merita più attenzione delle altre alternative per una ragione precisa: il
DIPME-DEVICE ne monta già uno. La domanda ``perché mmWave e non UWB?'' è quindi
attesa in sede di discussione, e la risposta va preparata.

L'UWB trasmette impulsi molto brevi su una banda larghissima (tipicamente
3--10~GHz), ottenendo un'eccellente risoluzione temporale e quindi spaziale, buona
penetrazione dei materiali e consumi molto contenuti (sotto i 10~mW in
trasmissione)~\cite{rfwireless}. È in grado di rilevare corpi fermi e, in
letteratura, anche parametri vitali.

Gli argomenti a favore del mmWave nel contesto specifico di questa tesi sono
quattro:

\begin{enumerate}[leftmargin=1.8em,itemsep=2pt]
  \item \textbf{maturità dei moduli consumer}: LD2410B e LD2420 costano 5--15~\euro,
        hanno firmware che espone direttamente presenza, distanza ed energia
        per-gate e non richiedono elaborazione a valle; i moduli UWB comparabili
        costano 20--50~\euro\ e tipicamente espongono dati grezzi che richiedono un
        processore vero;
  \item \textbf{dati per-gate già elaborati}: i venti canali di energia del
        LD2410B sono ciò che rende possibile l'indice di vitalità del
        \cref{cap:vitalita} su un microcontrollore da pochi euro;
  \item \textbf{destinazione d'uso}: l'UWB sul DIPME è pensato principalmente per
        localizzazione/ranging fra nodi, un problema diverso dal rilevamento di
        presenza in un volume;
  \item \textbf{complementarità, non concorrenza}: nulla vieta che il nodo finale
        usi entrambi. La raccomandazione del \cref{cap:conclusioni} è
        un'architettura a più sensori, non una sostituzione.
\end{enumerate}

\subsection{Ultrasuoni, CO\textsubscript{2}, visione}

Per completezza si citano tre alternative escluse a priori.
Gli \textbf{ultrasuoni} misurano distanza in modo economico ma sono molto
sensibili a tessuti fonoassorbenti e alla geometria della scena, e sotto un arredo
crollato la propagazione acustica è imprevedibile.
La \textbf{CO\textsubscript{2}} è un ottimo indicatore di occupazione di un
ambiente --- ed è già presente sul DIPME --- ma ha costanti di tempo dell'ordine
dei minuti e non localizza: inadatta a rispondere alla domanda ``c'è qualcuno
sotto \emph{questo} banco?''.
La \textbf{visione} (telecamere, termocamere) è la più informativa ma è esclusa da
vincoli di privacy in ambiente scolastico, dal costo e dall'inutilizzabilità nel
buio e nella polvere post-crollo.

\section{Posizionamento comparativo}
\label{sec:posizionamento}

La \cref{tab:comparativa} riassume il confronto fra le tre tecnologie candidate. I
valori sono da datasheet o da letteratura; la colonna mmWave è quella verificata
sperimentalmente nel \cref{cap:confronto}.

\begin{table*}[htbp]
\centering
\caption{Confronto fra le tecnologie di rilevamento di presenza candidate. Le
prestazioni della colonna mmWave sono verificate sperimentalmente nel
\cref{cap:confronto}.}
\label{tab:comparativa}
\small
\begin{tabularx}{\linewidth}{@{}lXXX@{}}
\toprule
\textbf{Parametro} & \textbf{PIR} & \textbf{mmWave 24~GHz} & \textbf{UWB 3--10~GHz} \\
\midrule
Rilevamento corpo fermo & no (fisicamente) & sì (micro-movimenti) & sì \\
Misura di distanza & no & sì & sì \\
Conteggio persone & no & limitato (no angolazione) & limitato \\
Rilevamento respiro & no & sì, parziale & sì \\
Penetrazione ostacoli & no & parziale (non metallo/cemento) & buona \\
Corrente tipica & $\sim$0,05~mA & 50--80~mA & $<$10~mW \\
Privacy & assenza di dato personale & alta (nessuna immagine) & alta \\
Costo indicativo & 1--3~\euro & 5--15~\euro & 20--50~\euro \\
Alimentazione & batteria adeguata & fissa preferibile & batteria adeguata \\
\bottomrule
\end{tabularx}
\end{table*}

Il quadro che emerge è quello di una \textbf{complementarità}, non di una
sostituzione: il PIR è imbattibile come guardiano a costo energetico nullo, il
mmWave è l'unico dei tre a offrire, a costo contenuto e con dati già elaborati,
una misura continua di \emph{quanto} si muove ciò che riflette. È esattamente la
grandezza necessaria all'indice di vitalità.

\medskip
\noindent
Come riferimento di prestazioni attese in condizioni ottimali, la letteratura
tecnica sui moduli mmWave consumer indica latenza di rilevamento del movimento
inferiore a 500~ms, conferma di persona ferma in 2--5~s, conferma di stanza vuota
in 20--60~s e tassi di falsi positivi inferiori a 0,05 eventi/h.
Il \cref{cap:confronto} verifica questi ordini di grandezza sul nostro esemplare.
